\section{Kesimpulan dan Saran}
\label{sec:kesimpulan}

\subsection{Kesimpulan Utama}
Penelitian dan kajian komprehensif ini berhasil mengintegrasikan teori fisika fundamental, penurunan analitik, kaidah seleksi teori grup, formalisme kuantum resonansi Kramers--Heisenberg--Dirac (KHD), instrumentasi nyata, analisis spektrum eksperimental, dan simulasi fisika maju pada spektroskopi Raman nanomaterial karbon berdimensi rendah serta heterostruktur 2D. Beberapa kesimpulan ilmiah utama yang diperoleh adalah sebagai berikut:

\begin{enumerate}
    \item \textbf{Penurunan Klasik dan Kuantum Hamburan Raman}: Modulasi polarizabilitas molekular oleh koordinat normal vibrasi kisi berhasil diturunkan secara eksak. Hamburan Raman fundamental mensyaratkan laju perubahan polarizabilitas tak-nol pada kesetimbangan ($(\partial\alpha/\partial Q)_0 \neq 0$). Formulasi kuantum menggunakan operator tangga osilator harmonik membuktikan kaidah seleksi transisi vibrasi fundamental $\Delta v = \pm 1$. Rasio intensitas anti-Stokes terhadap Stokes terbukti dikontrol secara simultan oleh faktor frekuensi kuartik dan hukum distribusi Boltzmann termal, yang memvalidasi penerapan spektroskopi Raman sebagai termometer kisi non-kontak.
    \item \textbf{Kaidah Seleksi dan Teori Grup Molekular}: Analisis simetri membuktikan Aturan Eksklusi Mutual pada sistem sentrosimetrik ($D_{\infty h}$ pada $\text{CO}_2$ dan $D_{6h}$ pada grafena), di mana mode dengan paritas genap ($g$) hanya aktif Raman dan inaktif IR. Pada grafena monolayer, mode optik in-plane simetri $E_{2g}$ menghasilkan pita $G$ pada $\approx 1582.4\text{ cm}^{-1}$. Ketiadaan pusat inversi pada kristal h-BN ($D_{3h}$) memungkinkan mode $E_{2g}$ ($E'$) aktif simultan pada Raman dan IR pada $\approx 1366.8\text{ cm}^{-1}$.
    \item \textbf{Teori Resonansi KHD dan Resonansi Ganda}: Formalisme gangguan kuantum KHD menjelaskan penguatan penampang lintang Raman hingga $10^3 - 10^6$ kali ketika energi foton laser cocok dengan transisi elektronik nyata. Pada grafena, struktur pita kerucut Dirac linier menghasilkan proses resonansi ganda (\textit{double resonance}) inter-valley dua fonon yang menjelaskan keaktifan pita $2D$ tanpa defek, sedangkan pita $D$ memerlukan defek untuk memenuhi kekekalan momentum kristal. Laju dispersi resonansi ganda terhadap energi foton terkonfirmasi bernilai $\partial\omega_{2D}/\partial E_L \approx 100\text{ cm}^{-1}/\text{eV}$.
    \item \textbf{Analisis Eksperimental Kritis Nanomaterial}:
    \begin{itemize}
        \item Grafena monolayer ditandai oleh pita $2D$ Lorentzian simetris tunggal dengan rasio $I_{2D}/I_G = 2.75$ dan FWHM $= 26.2\text{ cm}^{-1}$.
        \item Grafena \textit{bilayer} AB-Bernal terdekonvolusi menjadi empat komponen sub-pita Lorentzian ($2D_{11}$, $2D_{12}$, $2D_{21}$, $2D_{22}$) dengan FWHM gabungan $\approx 54\text{ cm}^{-1}$ dan rasio intensitas $I_{2D}/I_G \approx 1.03$.
        \item Grafena terdefek memperlihatkan pita $D$ ($1348.0\text{ cm}^{-1}$) dan $D'$ ($1621.5\text{ cm}^{-1}$) dengan $I_D/I_G = 0.90$, berkorespondensi dengan kerapatan defek titik $n_D \approx 2.02 \times 10^{11}\text{ cm}^{-2}$.
        \item SWCNT metalik membuktikan keberadaan kopling fonon-plasmon melalui profil asimetris Breit--Wigner--Fano (BWF) pada pita $G^-$ dengan parameter asimetri $1/q = -0.22$.
        \item Heterostruktur grafena/h-BN menunjukkan pemisahan tegas antara mode fonon substrat h-BN $E_{2g}$ ($1366.8\text{ cm}^{-1}$) dari pita defek grafena, serta penyempitan luar biasa FWHM pita $2D$ grafena menjadi $17.5\text{ cm}^{-1}$ akibat supresi fluktuasi muatan lokal substrat.
    \end{itemize}
    \item \textbf{Validasi Kuantitatif Pemodelan Fisika Maju}: Model fisika maju yang dibangun dari konstanta pegas ikatan analitik ($k_{\mathrm{C=C}} = 885.5\text{ N/m}$, $k_{\mathrm{B-N}} = 671.0\text{ N/m}$) dan tensor polarizabilitas menghasilkan kesesuaian tinggi dengan data eksperimen ($R^2 \ge 0.985$ dan korelasi Pearson $r > 0.99$ pada seluruh sampel). Deviasi frekuensi puncak maksimum sangat kecil ($\Delta\tilde{\nu} \le 0.8\text{ cm}^{-1}$), membuktikan keandalan pendekatan forward modeling yang diajukan.
\end{enumerate}

\subsection{Keterbatasan Pemodelan}
Beberapa keterbatasan model yang perlu digarisbawahi secara ilmiah meliputi:
\begin{enumerate}
    \item Pendekatan osilator harmonik tidak memperhitungkan pergeseran anharmonik orde tinggi (interaksi fonon-fonon cubic dan quartic) yang menyebabkan peluruhan intrinsik garis spektral pada temperatur tinggi.
    \item Model forward SWCNT mengasumsikan tabung terisolasi sempurna berorientasi seragam, sehingga belum memperhitungkan interaksi van der Waals antar-tabung dalam berkas (\textit{bundle effect}) yang dapat memodifikasi frekuensi mode RBM sebesar $5 - 15\text{ cm}^{-1}$.
    \item Efek dielektrik substrat h-BN dimodelkan melalui renormalisasi lebar garis dan pergeseran frekuensi efektif, bukan perhitungan struktur pita elektronik \textit{ab initio} yang menyertakan pola superkisi Moiré pada sudut puntir (\textit{twist angle}) tertentu.
\end{enumerate}

\subsection{Saran dan Pengembangan Riset Lanjutan}
Untuk memperluas cakupan dan kedalaman analisis pada masa mendatang, disarankan:
\begin{enumerate}
    \item Mengintegrasikan perhitungan kimia komputasi \textit{Density Functional Theory} (DFT) dan \textit{Density Functional Perturbation Theory} (DFPT) untuk menurunkan kurva dispersi fonon lengkap di seluruh zona Brillouin, tensor polarizabilitas dinamis Raman, serta intensitas kuantum Raman tanpa ketergantungan pada parameter empiris.
    \item Menerapkan pemetaan Raman konfokal beresolusi spasial tinggi (\textit{hyperspectral Raman imaging}) pada heterostruktur grafena/h-BN untuk memetakan distribusi regangan biaksial versus kerapatan muatan doping secara spasial melalui plot korelasi $\omega_{2D}$ versus $\omega_G$.
    \item Mengembangkan analisis \textit{in situ} ketergantungan potensial gerbang (\textit{gate voltage modulation}) dan temperatur kriogenik ($4.2\text{ K}$) untuk mengamati renormalisasi kopling elektron-fonon pada anomali Kohn secara langsung.
\end{enumerate}
